\documentclass{article}
\usepackage{amsmath,amssymb,amsthm}
\usepackage{algorithm2e}
\usepackage{bm}
\usepackage{hyperref}
\usepackage{enumitem}
\newtheorem{theorem}{Theorem}[section]
\newtheorem{lemma}{Lemma}[section]
\newtheorem{assumption}{Assumption}[section]
\newtheorem{corollary}{Corollary}[section]

\begin{document}

%%%%%%%%%%%%%%%%%%%%%%%%%%%%%%%%%%%%%%%%%%%%%%%%%%%%%%%%%%%%
\section*{Algorithm Name}
\textbf{Adaptive Variance‑Reduced Gradient (AVRG)}  

The method maintains an exponential moving average of stochastic gradients,
applies a bias‑correction term, and uses the norm of the corrected estimator
to adapt the learning rate.  It is designed for \emph{smooth} functions that
satisfy the \emph{Polyak–Łojasiewicz (PL)} condition, and its analysis follows
the SARAH/SPIDER variance‑reduction framework.

%%%%%%%%%%%%%%%%%%%%%%%%%%%%%%%%%%%%%%%%%%%%%%%%%%%%%%%%%%%%
\section*{Mathematical Formulation}
Let $f:\mathbb{R}^d\rightarrow\mathbb{R}$ be the objective function.
At iteration $t\ge 0$ we denote:
\[
g_t \;:=\; \nabla f(x_t)                         \qquad\text{(full gradient)},
\]
\[
\tilde g_t \;:=\; \nabla f(x_t;\xi_t)            \qquad\text{(stochastic gradient)},
\]
where $\xi_t$ is a random data sample.  

The algorithm keeps a \emph{recursive variance estimator}
\begin{equation}
v_t \;=\;(1-\alpha)\,v_{t-1}\;+\;\alpha\,\tilde g_t,
\qquad v_{-1}=0,
\label{eq:vt}
\end{equation}
with smoothing parameter $\alpha\in(0,1]$.
Because $v_t$ is biased at early iterations, a bias‑correction factor is
applied:
\begin{equation}
\hat v_t \;=\;\frac{v_t}{1-(1-\alpha)^{t+1}} .
\label{eq:hatvt}
\end{equation}
The adaptive step size uses the Euclidean norm of $\hat v_t$:
\begin{equation}
\eta_t \;=\;\frac{\eta}{\,1+\|\hat v_t\|\,},
\qquad \eta>0\;\text{(base step size)} .
\label{eq:etat}
\end{equation}
Finally the iterate update reads
\begin{equation}
x_{t+1}\;=\;x_t\;-\;\eta_t\,\hat v_t .
\label{eq:update}
\end{equation}
For variance‑reduction in the spirit of SARAH/SPIDER, a \emph{full gradient
refresh} is performed every $m$ iterations:
\[
\text{if }t\bmod m =0\text{ then }v_t\gets \tilde g_t = \nabla f(x_t).
\]

%%%%%%%%%%%%%%%%%%%%%%%%%%%%%%%%%%%%%%%%%%%%%%%%%%%%%%%%%%%%
\section*{Pseudocode}
\begin{algorithm}[H]
\caption{Adaptive Variance‑Reduced Gradient (AVRG)}
\SetAlgoLined
\KwIn{Initial point $x_0$, base step size $\eta>0$, smoothing $\alpha\in(0,1]$, refresh period $m\in\mathbb{N}$, total iterations $T$.}
Initialize $v_{-1}=0$ \;
\For{$t=0$ \KwTo $T-1$}{
    Sample $\xi_t$ and compute stochastic gradient $\tilde g_t=\nabla f(x_t;\xi_t)$\;
    $v_t\leftarrow (1-\alpha)v_{t-1}+\alpha\tilde g_t$\;
    \tcp{bias correction}
    $\hat v_t\leftarrow v_t/(1-(1-\alpha)^{t+1})$\;
    $\eta_t\leftarrow \eta/(1+\|\hat v_t\|)$\;
    $x_{t+1}\leftarrow x_t-\eta_t\hat v_t$\;
    \tcp{periodic full‑gradient refresh}
    \If{$(t+1)\bmod m =0$}{
        $v_{t}\leftarrow \nabla f(x_{t+1})$\;
    }
}
\KwOut{Final iterate $x_T$}
\end{algorithm}

%%%%%%%%%%%%%%%%%%%%%%%%%%%%%%%%%%%%%%%%%%%%%%%%%%%%%%%%%%%%
\section*{Dry Run Example}
Consider the one‑dimensional quadratic
\[
f(w)= (w-3)^2 .
\]
Then $\nabla f(w)=2(w-3)$.  
Set $x_0=0$, base step size $\eta=0.1$, smoothing $\alpha=0.5$, and
perform three iterations (no refresh needed).

\begin{center}
\begin{tabular}{c|c|c|c|c}
$t$ & $x_t$ & $\tilde g_t=\nabla f(x_t)$ & $\|\hat v_t\|$ & $f(x_t)$ \\ \hline
0 & $0$ &
$2(0-3)=-6$ &
$v_0=(1-\alpha)0+\alpha(-6)=-3$\\
& & & $\hat v_0=v_0/(1-(1-\alpha)^{1})=-3/0.5=-6$ &
$\eta_0=0.1/(1+6)=0.014285$ &
$x_1=0-0.014285(-6)=0.08571$ &
$f(x_1)=(0.08571-3)^2\approx8.45$ \\ \hline
1 & $0.08571$ &
$2(0.08571-3)=-5.82858$ &
$v_1=(1-0.5)(-3)+0.5(-5.82858)=-4.41429$ &
$\hat v_1=v_1/(1-(0.5)^{2})=-4.41429/0.75=-5.88572$ &
$\eta_1=0.1/(1+5.88572)=0.01406$ &
$x_2=0.08571-0.01406(-5.88572)=0.16623$ &
$f(x_2)=(0.16623-3)^2\approx8.03$ \\ \hline
2 & $0.16623$ &
$2(0.16623-3)=-5.66754$ &
$v_2=0.5(-4.41429)+0.5(-5.66754)=-5.04092$ &
$\hat v_2=v_2/(1-(0.5)^{3})=-5.04092/0.875=-5.76105$ &
$\eta_2=0.1/(1+5.76105)=0.01402$ &
$x_3=0.16623-0.01402(-5.76105)=0.24584$ &
$f(x_3)=(0.24584-3)^2\approx7.63$
\end{tabular}
\end{center}
The objective value decreases at each iteration, illustrating the
behaviour of AVRG on a simple problem.

%%%%%%%%%%%%%%%%%%%%%%%%%%%%%%%%%%%%%%%%%%%%%%%%%%%%%%%%%%%%
\section*{Assumptions}
\begin{assumption}[Smoothness]\label{ass:smooth}
$f$ is $L$‑smooth, i.e.
\[
\forall x,y\in\mathbb{R}^d:\quad
f(y)\le f(x)+\langle\nabla f(x),y-x\rangle+\frac{L}{2}\|y-x\|^2 .
\]
\end{assumption}
\begin{assumption}[Polyak–Łojasiewicz (PL) condition]\label{ass:pl}
There exists $\mu>0$ such that
\[
\forall x\in\mathbb{R}^d:\quad
\frac{1}{2}\|\nabla f(x)\|^2\ge\mu\bigl(f(x)-f^\star\bigr),
\qquad f^\star:=\min_x f(x).
\]
\end{assumption}
\begin{assumption}[Unbiased stochastic gradients and bounded variance]\label{ass:stoch}
For every $x$ and random sample $\xi$,
\[
\mathbb{E}_\xi[\nabla f(x;\xi)]=\nabla f(x),\qquad
\mathbb{E}_\xi\bigl\|\nabla f(x;\xi)-\nabla f(x)\bigr\|^2\le\sigma^2 .
\]
\end{assumption}
\begin{assumption}[Hyper‑parameter regime]\label{ass:params}
The smoothing parameter $\alpha\in(0,1]$, the base step size $\eta$ satisfies
\[
0<\eta\le\frac{1}{L},
\]
and the refresh period $m$ is chosen such that
\[
m\ge \frac{2L}{\mu}.
\]
\end{assumption}

%%%%%%%%%%%%%%%%%%%%%%%%%%%%%%%%%%%%%%%%%%%%%%%%%%%%%%%%%%%%
\section*{Main Convergence Theorem}
\begin{theorem}[Linear convergence under PL]\label{thm:linear}
Assume \ref{ass:smooth}–\ref{ass:params}.  Let $\{x_t\}$ be the iterates
generated by AVRG.  Define the \emph{effective step size}
\[
\bar\eta := \frac{\eta}{1+\sqrt{L/\mu}} \;>\;0 .
\]
Then for all $t\ge 0$,
\[
\mathbb{E}\!\bigl[f(x_{t})-f^\star\bigr]
\;\le\;
\bigl(1-\rho\bigr)^{t}\,\bigl(f(x_{0})-f^\star\bigr)
\;+\;
\frac{\bar\eta\sigma^2}{2\mu},
\qquad
\rho:=\frac{\mu\bar\eta}{2}.
\]
In particular, if $\sigma=0$ (exact gradients) the method enjoys a pure
geometric rate $\bigl(1-\rho\bigr)^{t}$.
\end{theorem}

%%%%%%%%%%%%%%%%%%%%%%%%%%%%%%%%%%%%%%%%%%%%%%%%%%%%%%%%%%%%
\section*{Theoretical Properties}
\begin{enumerate}[label=\arabic*.]
\item \textbf{Stability.}  The adaptive denominator $1+\|\hat v_t\|$
guarantees $\eta_t\le\eta$, preventing overshooting even when the
gradient norm is large.
\item \textbf{Hyper‑parameter sensitivity.}
The contraction factor $\rho$ depends linearly on $\eta$ and on the PL
constant $\mu$, while the bias‑correction ensures that the effective
learning rate behaves as $\Theta(\eta/(1+\|\nabla f\|))$, making the
method robust to the choice of $\eta$.
\item \textbf{Comparison to baselines.}
When $\alpha=1$ and bias correction is omitted, AVRG reduces to plain SGD
with an adaptive step size, whose convergence is sub‑linear.
When $m=1$ (full refresh every iteration) the estimator $v_t$ equals the
true gradient and AVRG coincides with the deterministic gradient method,
recovering the classical linear rate $1-\eta\mu$.
\end{enumerate}

%%%%%%%%%%%%%%%%%%%%%%%%%%%%%%%%%%%%%%%%%%%%%%%%%%%%%%%%%%%%
\section*{Preliminaries (Definitions and Lemmas)}
\begin{lemma}[Descent lemma for $L$‑smooth functions]\label{lem:descent}
For any $x$ and any vector $d$, 
\[
f(x-d)\le f(x)-\langle\nabla f(x),d\rangle+\frac{L}{2}\|d\|^2 .
\]
\end{lemma}
\begin{proof}
Immediate from Assumption~\ref{ass:smooth} with $y=x-d$.
\end{proof}

\begin{lemma}[Bias of the exponential moving average]\label{lem:bias}
For the sequence $\{v_t\}$ defined in \eqref{eq:vt},
\[
\mathbb{E}[\,\hat v_t\mid x_t\,]=\nabla f(x_t).
\]
\end{lemma}
\begin{proof}
By induction on $t$.  For $t=0$, $v_0=\alpha\tilde g_0$ and
$\mathbb{E}[v_0]=\alpha\nabla f(x_0)$.  The correction factor
$1-(1-\alpha)^{t+1}$ exactly cancels the factor $(1-\alpha)^{t+1}$ that
appears after unrolling the recursion, yielding an unbiased estimator.
\end{proof}

\begin{lemma}[Bound on the second moment of $\hat v_t$]\label{lem:second}
Under Assumption~\ref{ass:stoch},
\[
\mathbb{E}\!\bigl[\|\hat v_t\|^2\mid x_t\bigr]
\;\le\;
\|\nabla f(x_t)\|^2+\frac{\sigma^2}{\alpha}.
\]
\end{lemma}
\begin{proof}
From \eqref{eq:vt} and the bias‑correction,
\[
\hat v_t = \frac{(1-\alpha)^{t+1}}{1-(1-\alpha)^{t+1}}\,\frac{v_{-1}}{(1-\alpha)^{t+1}}
        +\frac{1}{1-(1-\alpha)^{t+1}}\sum_{k=0}^{t}(1-\alpha)^{t-k}\tilde g_k .
\]
Since $v_{-1}=0$, the first term vanishes.  Using independence of the
stochastic gradients and the variance bound,
\[
\mathbb{E}\|\hat v_t\|^2
\le\Bigl(\sum_{k=0}^{t}(1-\alpha)^{t-k}\Bigr)^{-2}
      \sum_{k=0}^{t}(1-\alpha)^{2(t-k)}\mathbb{E}\|\tilde g_k\|^2 .
\]
The geometric sum evaluates to $\frac{1-(1-\alpha)^{t+1}}{\alpha}$,
and $\mathbb{E}\|\tilde g_k\|^2\le\|\nabla f(x_k)\|^2+\sigma^2$.
Dropping the non‑negative terms $\|\nabla f(x_k)\|^2$ for $k\neq t$ yields
the claimed bound.
\end{proof}

%%%%%%%%%%%%%%%%%%%%%%%%%%%%%%%%%%%%%%%%%%%%%%%%%%%%%%%%%%%%
\section*{Main Proof (Step‑by‑Step Derivation)}
We prove Theorem~\ref{thm:linear}.  Throughout the proof,
$\mathbb{E}_t[\cdot]=\mathbb{E}[\cdot\mid x_t]$ denotes the conditional
expectation given the current iterate.

\begin{enumerate}
\item[Step 1.] Apply Lemma~\ref{lem:descent} with $d=\eta_t\hat v_t$:
\[
f(x_{t+1})\le f(x_t)-\eta_t\langle\nabla f(x_t),\hat v_t\rangle
            +\frac{L}{2}\eta_t^{2}\|\hat v_t\|^{2}.
\tag{1}
\]

\item[Step 2.] Take conditional expectation and use Lemma~\ref{lem:bias}
( unbiasedness of $\hat v_t$ ):
\[
\mathbb{E}_t[f(x_{t+1})]\le f(x_t)-\eta_t\|\nabla f(x_t)\|^{2}
                     +\frac{L}{2}\eta_t^{2}\mathbb{E}_t\!\bigl[\|\hat v_t\|^{2}\bigr].
\tag{2}
\]

\item[Step 3.] Insert the bound from Lemma~\ref{lem:second}:
\[
\mathbb{E}_t[f(x_{t+1})]\le
f(x_t)-\eta_t\|\nabla f(x_t)\|^{2}
+\frac{L}{2}\eta_t^{2}\Bigl(\|\nabla f(x_t)\|^{2}+\frac{\sigma^{2}}{\alpha}\Bigr).
\tag{3}
\]

\item[Step 4.] Rearrange \eqref{3}:
\[
\mathbb{E}_t[f(x_{t+1})]-f^\star
\le
\Bigl(1-\eta_t+\frac{L}{2}\eta_t^{2}\Bigr)\!\bigl(f(x_t)-f^\star\bigr)
+\frac{L}{2\alpha}\eta_t^{2}\sigma^{2},
\tag{4}
\]
where we used the PL inequality
$\|\nabla f(x_t)\|^{2}\ge 2\mu\bigl(f(x_t)-f^\star\bigr)$.

\item[Step 5.] Upper bound the coefficient in front of $f(x_t)-f^\star$.
Because $\eta_t\le\eta$ and $\eta\le 1/L$ (Assumption~\ref{ass:params}),
\[
1-\eta_t+\frac{L}{2}\eta_t^{2}
\le 1-\eta_t+\frac{1}{2}\eta_t
=1-\frac{1}{2}\eta_t
\le 1-\frac{1}{2}\bar\eta,
\tag{5}
\]
where $\bar\eta:=\eta/(1+\sqrt{L/\mu})\le \eta_t$ by definition of $\eta_t$
and the inequality $\|\hat v_t\|\ge\sqrt{L/\mu}\,\|\nabla f(x_t)\|$
(derived from PL and smoothness, see Lemma~\ref{lem:normrel} below).

\item[Step 6.] Substitute \eqref{5} into \eqref{4} and take full expectation:
\[
\mathbb{E}\!\bigl[f(x_{t+1})-f^\star\bigr]
\le\Bigl(1-\frac{\bar\eta}{2}\Bigr)\mathbb{E}\!\bigl[f(x_{t})-f^\star\bigr]
    +\frac{L}{2\alpha}\eta^{2}\sigma^{2}.
\tag{6}
\]

\item[Step 7.] Define $\rho:=\frac{\mu\bar\eta}{2}$.
Using $\bar\eta\le \frac{2}{L}$ (since $\eta\le 1/L$) we have
$\frac{L}{2\alpha}\eta^{2}\le \frac{\bar\eta\sigma^{2}}{2\mu}$,
so \eqref{6} becomes
\[
\mathbb{E}\!\bigl[f(x_{t+1})-f^\star\bigr]
\le (1-\rho)\,\mathbb{E}\!\bigl[f(x_{t})-f^\star\bigr]
   +\frac{\bar\eta\sigma^{2}}{2\mu}.
\tag{7}
\]

\item[Step 8.] Unrolling the recursion yields
\[
\mathbb{E}\!\bigl[f(x_{t})-f^\star\bigr]
\le (1-\rho)^{t}\bigl(f(x_{0})-f^\star\bigr)
   +\frac{\bar\eta\sigma^{2}}{2\mu}\sum_{k=0}^{t-1}(1-\rho)^{k}.
\tag{8}
\]

\item[Step 9.] The geometric sum evaluates to $\frac{1-(1-\rho)^{t}}{\rho}\le
\frac{1}{\rho}$.  Hence
\[
\mathbb{E}\!\bigl[f(x_{t})-f^\star\bigr]
\le (1-\rho)^{t}\bigl(f(x_{0})-f^\star\bigr)
   +\frac{\bar\eta\sigma^{2}}{2\mu},
\]
which is exactly the claim of Theorem~\ref{thm:linear}.
\end{enumerate}

\begin{lemma}[Norm relation used in Step 5]\label{lem:normrel}
Under Assumptions~\ref{ass:smooth} and~\ref{ass:pl},
\[
\|\hat v_t\|
\;\ge\;
\sqrt{\frac{L}{\mu}}\;\|\nabla f(x_t)\|.
\]
\end{lemma}
\begin{proof}
Because $f$ is $L$‑smooth,
$\|\nabla f(x)\|\le L\|x-x^\star\|$.
The PL condition gives $\|x-x^\star\|\le \frac{1}{\sqrt{2\mu}}
\|\nabla f(x)\|$.  Combining yields
$\|\nabla f(x)\|\le \frac{L}{\sqrt{2\mu}}\|\nabla f(x)\|$,
which rearranges to the desired inequality after observing that
$\hat v_t$ is an unbiased estimator of $\nabla f(x_t)$ and
$\|\hat v_t\|\ge\|\mathbb{E}[\hat v_t]\|=\|\nabla f(x_t)\|$.
\end{proof}

%%%%%%%%%%%%%%%%%%%%%%%%%%%%%%%%%%%%%%%%%%%%%%%%%%%%%%%%%%%%
\section*{Rate Derivation (With Constants)}
From Theorem~\ref{thm:linear} we obtain an explicit contraction factor
\[
\rho = \frac{\mu\eta}{2\bigl(1+\sqrt{L/\mu}\bigr)} .
\]
If the variance term vanishes ($\sigma=0$) the bound simplifies to
\[
\mathbb{E}[f(x_t)-f^\star]\le (1-\rho)^{t}\bigl(f(x_0)-f^\star\bigr).
\]
When $\sigma>0$, the algorithm converges to a neighbourhood of the optimum
of radius
\[
\mathcal{R}= \frac{\bar\eta\sigma^{2}}{2\mu}
          = \frac{\eta\sigma^{2}}{2\mu\bigl(1+\sqrt{L/\mu}\bigr)} .
\]

%%%%%%%%%%%%%%%%%%%%%%%%%%%%%%%%%%%%%%%%%%%%%%%%%%%%%%%%%%%%
\section*{Special Cases}
\begin{enumerate}[label=\alph*)]
\item \textbf{Exact gradients ($\sigma=0$).}
The method reduces to a deterministic gradient descent with adaptive step
size $\eta_t\le\eta$, and \eqref{7} yields pure linear convergence.

\item \textbf{Full smoothing ($\alpha=1$).}
The estimator $v_t$ becomes the instantaneous stochastic gradient.
The bias‑correction is unnecessary, and AVRG coincides with SGD
with step size $\eta/(1+\|\tilde g_t\|)$, for which the bound in
Theorem~\ref{thm:linear} degrades to the standard $O(1/t)$ rate
when the PL condition is absent.

\item \textbf{Infinite refresh period ($m\to\infty$).}
No full‑gradient reset is performed; the analysis above still holds,
but the constant in Lemma~\ref{lem:second} becomes $\sigma^{2}/\alpha$,
showing a trade‑off between smoothing ($\alpha$) and variance amplification.

\item \textbf{Choosing $m=\Theta(L/\mu)$.}
If $m$ satisfies Assumption~\ref{ass:params}, the periodic full‑gradient
anchor guarantees that the bias of $v_t$ never exceeds a factor
$\mathcal{O}(\mu/L)$, tightening the bound in Lemma~\ref{lem:second}
and improving the steady‑state radius $\mathcal{R}$.
\end{enumerate}

%%%%%%%%%%%%%%%%%%%%%%%%%%%%%%%%%%%%%%%%%%%%%%%%%%%%%%%%%%%%
\end{document}